\section{The Lubotzky--Wilson embedding}

Fix a prime $p$. This section proves the consequence of
\cite{LubotzkyWilson1984} needed for Camina's theorem: each countably
based pro-$p$ group embeds in a topologically finitely generated
pro-$p$ group. The argument follows the source's wreath-product
construction, supplies the elementary arithmetic proof omitted there,
and makes continuity and closedness explicit. For finite $p$-groups,
the first lemma supplies directly the derived-subgroup embedding for
which that paper cites Neumann and Neumann. We use
$[a,b]=aba^{-1}b^{-1}$ throughout this section.

\begin{lemma}[Embedding a finite group in a derived subgroup]
\label{lem:finite-p-derived-embedding}
\source{lubotzky-wilson1984:page-0002}
\lean{Camina1997.constantWreath_mem_commutator, Camina1997.finiteDerivedEmbedding, Camina1997.finiteDerivedEmbedding_injective, Camina1997.exists_finite_pGroup_derived_embedding}
\leanok
Every finite $p$-group $A$ embeds in $F'$ for some finite $p$-group $F$.
\end{lemma}

\begin{proof}
Put $n=|A|$ and let $C_n=\mathbf Z/n\mathbf Z$, written additively.
Let $B$ be the group of functions $C_n\to A$ with pointwise
multiplication. Let $C_n$ act on $B$ by
$(q\cdot f)(x)=f(x-q)$, and put $F=B\rtimes C_n$.
Its order is $n^n n$, a power of $p$; this also covers $n=1$.
For $a\in A$, let $d(a)\in B$ be the constant function with value $a$.
The map $d$ is an injective homomorphism. Define $f_a(\bar k)=a^k$.
This is well-defined because $a^n=1$ by Lagrange's theorem.
Let $t$ denote the element $-1\in C_n$ in the acting subgroup
of $F$. Conjugation by $t$ sends $f(x)$ to $f(x+1)$, so for every $k$,
\[
 [t,f_a](\bar k)=f_a(\overline{k+1})f_a(\bar k)^{-1}
             =a^{k+1}a^{-k}=a.
\]
Thus $d(a)=[t,f_a]\in F'$, as required. The functions $f_a$
need not depend homomorphically on $a$; the embedding is the
constant-function map $d$.
This gives directly the finite $p$-group case of the
derived-subgroup embedding invoked in \cite{LubotzkyWilson1984}, p.~398.
\end{proof}

\begin{lemma}[Separation of sparse powers]
\label{lem:sparse-powers-separation}
\source{lubotzky-wilson1984:page-0002}
\lean{Camina1997.sparseAtom_sum_injective, Camina1997.sparse_pow_shift_eq_zero, Camina1997.sparse_pow_injective, Camina1997.sparse_powers_separation}
\leanok
Let $q,r,s,i,j$ be natural numbers with $q>2$, and let
$k\in\mathbf Z/q^{s+1}\mathbf Z$. Suppose that $k+q^i$ and
$k+q^j$ are distinct and each belongs to
$\{q^{r+1},\ldots,q^s\}$ in this cyclic group.
Then $r<i,j\le s$ and $k=0$.
\end{lemma}

\begin{proof}
Consider the sequence $a_0=0$ and $a_{u+1}=q^u$.
If $u<v$, then $2a_u<a_v$: for $u=0$ this is positivity,
and otherwise $a_v\ge qa_u>2a_u$.
Consequently two sums of two terms are equal only when their
unordered pairs of indices agree. To prove this, order the
indices in each pair as $u\le v$ and $u'\le v'$.
If $v<v'$, then $a_u+a_v\le2a_v<a_{v'}\le a_{u'}+a_{v'}$,
contradicting equality. Interchanging the pairs excludes $v'<v$,
so $v=v'$. Cancellation and strict increase then give $u=u'$.

Write $M=q^{s+1}$ and put $Q_e=q^e$ if $e\le s$, and $Q_e=0$
if $e>s$. In the latter case $q^e$ is divisible by $M$, so
$Q_e$ represents $q^e$ modulo $M$ in both cases.
Choose $a,b$ in the stated block with $k+q^i=q^a$ and
$k+q^j=q^b$. Distinctness implies $a\ne b$. Eliminating $k$
gives
\[
 Q_i+q^b\equiv q^a+Q_j\pmod M.
\]
Each side is a natural number at most $2q^s<M$, so the congruence
is equality in $\mathbf N$. The uniqueness of unordered pairs
just proved applies to these four terms of the sequence $(a_u)$.
The matching $q^b=q^a$ is impossible since $a\ne b$.
Thus $Q_i=q^a$ and $Q_j=q^b$. Both right sides are positive,
so $i,j\le s$; strict increase of the powers of $q$ gives
$i=a$ and $j=b$. Substitution in $k+q^i=q^a$ gives $k=0$,
and the bounds on $a,b$ give the conclusion.
\end{proof}

The printed arithmetic lemma assumes $r>0$. The proof above also
allows $r=0$ and zero values of the initial exponents.
It uses only $q>2$, without a primality assumption.

\begin{lemma}[Closed subgroups of products of finite $p$-groups]
\label{lem:closed-subgroup-finite-p-product}
\uses{def:pro-p-group}
\lean{Camina1997.IsProPGroup.pi_of_finite_discrete, Camina1997.IsProPGroup.of_isClosed_subgroup}
\leanok
If $(W_i)_{i\in I}$ is any family of finite discrete $p$-groups,
then $W=\prod_{i\in I}W_i$ is pro-$p$, and every closed subgroup
of $W$ is pro-$p$ with its subspace topology.
\end{lemma}

\begin{proof}
Each factor is compact Hausdorff, so their product is compact
Hausdorff by the product compactness theorem. Coordinatewise
multiplication and inversion are continuous, since their
composites with every coordinate projection are continuous.
A closed subgroup is compact Hausdorff with continuous group
operations inherited from $W$.
For a finite set of indices $S$, the subgroup
\[
 U_S=H\cap\bigcap_{n\in S}\ker(W\longrightarrow W_n)
\]
is open and normal in $H$. These subgroups form an identity
neighbourhood basis: a basic product neighbourhood restricts
finitely many coordinates, and each restricted neighbourhood
of the identity contains its singleton.
Projection identifies $H/U_S$ with a subgroup of
$\prod_{n\in S}W_n$. This product has $p$-power order; its
subgroups do too by Lagrange's theorem. Thus every $U_S$ has
finite $p$-power index, which is the definition of pro-$p$.
\end{proof}

\begin{lemma}[Countably many finite quotients separate points]
\label{lem:countable-pro-p-quotient-embedding}
\uses{lem:pro-p-countable-quotient-tower,lem:compact-injection-closed-embedding}
\source{lubotzky-wilson1984:page-0001,lubotzky-wilson1984:page-0002}
\lean{Camina1997.IsProPGroup.quotientDiagonal, Camina1997.IsProPGroup.isClosedEmbedding_quotientDiagonal, Camina1997.IsProPGroup.exists_closedEmbedding_quotient_product}
\leanok
Every countably based pro-$p$ group $P$ has a continuous injective
homomorphism into a countable product of finite discrete $p$-groups.
This map is a homeomorphism onto a closed subgroup.
\end{lemma}

\begin{proof}
Choose a decreasing identity neighbourhood basis $(N_n)$ of open
normal subgroups from Lemma~\ref{lem:pro-p-countable-quotient-tower}.
Put $A_n=P/N_n$, a finite $p$-group with its discrete topology.
The diagonal quotient map $e:P\to\prod_n A_n$ is a homomorphism.
Each coordinate is continuous because the inverse image of each
singleton is an open coset of $N_n$. Hence $e$ is continuous by
the definition of the product topology.
If $x\in\ker e$, then $x\in N_n$ for every $n$.
The intersection of this basis is trivial, so $x=1$.
Thus the kernel is trivial. The compact-source lemma applies,
since $P$ is compact and the product is Hausdorff.
\end{proof}

\begin{lemma}[Sparse coordinates for a sequence of finite groups]
\label{lem:wreath-sparse-coordinates}
\source{lubotzky-wilson1984:page-0002}
\lean{Camina1997.LubotzkyWilson.blockEnd_strictMono, Camina1997.LubotzkyWilson.powerCoordinate_injective, Camina1997.LubotzkyWilson.sparseFunction_elementIndex, Camina1997.LubotzkyWilson.sparseFunction_ne_one_imp}
\leanok
Let $(F_n)_{n\ge0}$ be finite groups and let $q>2$ be a natural
number. Put $r_0=0$ and $r_{n+1}=r_n+|F_n|$, and choose bijections
\[
 e_n:F_n\longrightarrow\{r_n+1,\ldots,r_{n+1}\}.
\]
The sequence $(r_n)$ is strictly increasing, and its indicated
blocks are pairwise disjoint. In
$C_n=\mathbf Z/q^{r_{n+1}+1}\mathbf Z$, the coordinates
$q^{e_n(a)}$ for $a\in F_n$ are distinct. Thus there is a function
$\theta_n:C_n\to F_n$ with
\[
 \theta_n(q^{e_n(a)})=a,
 \qquad
 \theta_n(k)=1\quad\text{off these coordinates}.
\]
In particular, $\theta_n(k)\ne1$ implies that
$k=q^i$ for some $r_n<i\le r_{n+1}$.
\end{lemma}

\begin{proof}
Every group contains its identity, so $|F_n|>0$ and
$r_{n+1}>r_n$. The interval $(r_n,r_{n+1}]$ has exactly
$|F_n|$ integer elements and hence admits the stated bijection.
If $m<n$, then $r_{m+1}\le r_n$, so these intervals are disjoint.
For $r_n<i\le r_{n+1}$, the natural number $q^i$ is strictly
less than $q^{r_{n+1}+1}$. Distinct exponents give distinct
natural powers since $q>1$, and these representatives remain
distinct modulo $q^{r_{n+1}+1}$. We may therefore prescribe
the indicated values of $\theta_n$ independently and set it
to $1$ elsewhere. The final support assertion follows directly.
\end{proof}

\begin{lemma}[Isolating commutators in the wreath products]
\label{lem:wreath-isolated-commutators}
\uses{lem:sparse-powers-separation,lem:wreath-sparse-coordinates}
\source{lubotzky-wilson1984:page-0002,lubotzky-wilson1984:page-0003}
\lean{Camina1997.LubotzkyWilson.shifted_commutator_support, Camina1997.LubotzkyWilson.cyclicShift_commutator_eq_single, Camina1997.LubotzkyWilson.cyclicShift_commutator_eq_one_of_not_mem, Camina1997.LubotzkyWilson.isolated_coordinate_commutator, Camina1997.LubotzkyWilson.coordinate_commutator_subgroup}
\leanok
Use the finite groups, cyclic coordinate groups, exponents, and
sparse functions of Lemma~\ref{lem:wreath-sparse-coordinates}.
Form the regular wreath products
\[
 W_n=F_n^{C_n}\rtimes C_n,
 \qquad(c\cdot f)(k)=f(k-c),
 \qquad W=\prod_{n\ge0}W_n.
\]
Let $t_n$ be the acting element $-1\in C_n$, and regard $\theta_n$
as an element of the base group of $W_n$. Put
$t=(t_n)_n$ and $\theta=(\theta_n)_n$.
Let $\delta_n:F_n\hookrightarrow W$ place its value at the zero
coordinate of the $n$th base group, with identity values elsewhere.
For all $a,b\in F_n$,
\[
 [t^{q^{e_n(a)}}\theta t^{-q^{e_n(a)}},
   t^{q^{e_n(b)}}\theta t^{-q^{e_n(b)}}]
   =\delta_n([a,b]).
\]
Consequently $\delta_n(F_n')\le\langle t,\theta\rangle$ for every $n$.
\end{lemma}

\begin{proof}
Conjugation by $t_m^u$, the acting element $-u$, sends a base
function $f(k)$ to $f(k+u)$. The commutator of two base functions
is their pointwise commutator. Hence, with $i=e_n(a)$ and $j=e_n(b)$,
the $m$th coordinate of the displayed commutator has value
\[
 [\theta_m(k+q^i),\theta_m(k+q^j)]
 \qquad(k\in C_m).
\]
If this is nontrivial, both values of $\theta_m$ are nonidentity
and the two arguments are distinct; equal arguments would give
the commutator of an element with itself. The support assertion
and the separation lemma give
\[
 r_m<i,j\le r_{m+1},\qquad k=0.
\]
The exponent $i$ lies in the $n$th block, so disjointness of the
blocks forces $m=n$. Thus all coordinates except the zero
coordinate of the $n$th factor are trivial. At that coordinate,
the prescribed sparse values give $[a,b]$. This proves the equality,
including the case $a=b$, when both sides are the identity.

The subgroup $\langle t,\theta\rangle$ contains all powers,
inverses, products, and commutators of its elements. The displayed
commutator therefore lies in it. Since $F_n'$ is generated by the
commutators $[a,b]$ and $\delta_n$ is a homomorphism, it contains
$\delta_n(F_n')$.
\end{proof}

\begin{lemma}[The product of derived groups in a two-generator group]
\label{lem:derived-product-two-generator}
\uses{lem:wreath-isolated-commutators,lem:closed-subgroup-finite-p-product,lem:compact-injection-closed-embedding,def:topologically-finitely-generated}
\source{lubotzky-wilson1984:page-0002,lubotzky-wilson1984:page-0003}
\lean{Camina1997.LubotzkyWilson.twoGeneratorHull, Camina1997.LubotzkyWilson.twoGenerated_twoGeneratorHull, Camina1997.LubotzkyWilson.topologicallyFG_twoGeneratorHull, Camina1997.LubotzkyWilson.isProP_twoGeneratorHull, Camina1997.LubotzkyWilson.derivedProductEmbedding, Camina1997.LubotzkyWilson.derivedProductEmbedding_isClosedEmbedding}
\leanok
Let $(F_n)_{n\ge0}$ be finite discrete $p$-groups, put $q=p^2$,
and form $W,t,\theta$ as in Lemma~\ref{lem:wreath-isolated-commutators}.
Then
\[
 H=\overline{\langle t,\theta\rangle}\le W
\]
is pro-$p$, and the two indicated elements generate a dense
subgroup of $H$. There is a closed group embedding
$\prod_nF_n'\longrightarrow H$ that places each entry at the
zero coordinate of its corresponding base group.
\end{lemma}

\begin{proof}
Since $p\ge2$, the number $q=p^2$ satisfies $q\ge4>2$,
so the preceding construction applies.
The closure of a subgroup of a topological group is a subgroup:
if $x,y$ are in the closure and $U$ is a neighbourhood of $xy^{-1}$,
continuity supplies neighbourhoods $V$ of $x$ and $Z$ of $y$
with $VZ^{-1}\subseteq U$; choose subgroup elements in $V$ and $Z$
to see that $U$ meets the subgroup. Thus $H$ is closed.
The subgroup generated by the two elements viewed in $H$
maps under the inclusion to exactly $\langle t,\theta\rangle$.
Its closure in the subspace $H$ is therefore $H$, proving
the assertion about the actual generating pair.
Each $W_n$ has $p$-power order: its base has order
$|F_n|^{q^{r_{n+1}+1}}$, and its cyclic quotient has order
$q^{r_{n+1}+1}=p^{2(r_{n+1}+1)}$.
The closed-subgroup lemma therefore makes $H$ pro-$p$.

Map $D=\prod_nF_n'$ into $W$ by placing each entry at the
zero coordinate of its corresponding base group and putting
$1$ at all other coordinates. This is a continuous injective
homomorphism, since it is one on each finite coordinate.
It is a closed embedding by compactness of $D$ and Hausdorffness
of $W$. Every finite-support element of its image is a finite
product of elements of the groups $\delta_n(F_n')$, and hence
belongs to $H$ by the previous lemma.
These finite-support elements are dense in its image: every
product neighbourhood specifies only finitely many factors,
and setting all other factors to $1$ gives such an element
in that neighbourhood. Since $H$ is closed, it contains the
whole image of $D$. Restricting the codomain to $H$ preserves
continuity and injectivity, and the compact-source lemma again
gives the required closed embedding.

The source chooses an odd prime or $4$ as its sparse-coordinate
base. The uniform choice $q=p^2$ has the two properties used in
that argument: $q>2$ and all cyclic factors have $p$-power order.
\end{proof}

\begin{theorem}[Lubotzky--Wilson for pro-$p$ groups]
\label{thm:lubotzky-wilson-pro-p}
\uses{lem:finite-p-derived-embedding,lem:countable-pro-p-quotient-embedding,lem:derived-product-two-generator,lem:compact-injection-closed-embedding}
\source{lubotzky-wilson1984:page-0001,lubotzky-wilson1984:page-0002,lubotzky-wilson1984:page-0003,camina1997:page-0012}
\lean{Camina1997.LubotzkyWilson.finiteDerivedProductHom, Camina1997.LubotzkyWilson.finiteDerivedProductHom_isClosedEmbedding, Camina1997.LubotzkyWilson.exists_topologicallyFG_proP_closedEmbedding}
\leanok
For every countably based pro-$p$ group $P$, there are a
topologically finitely generated pro-$p$ group $H$ and a
continuous injective homomorphism $P\to H$ that is a
homeomorphism onto a closed subgroup.
\end{theorem}

\begin{proof}
Choose the closed diagonal embedding
$P\to\prod_n A_n$, where the $A_n$ are finite discrete
$p$-group quotients, from
Lemma~\ref{lem:countable-pro-p-quotient-embedding}.
For each $n$, put
\[
 F_n=A_n\wr(\mathbf Z/|A_n|\mathbf Z).
\]
The first lemma makes $F_n$ a finite $p$-group and supplies
the injective constant-function homomorphism $e_n:A_n\to F_n'$.
Give $F_n$ the discrete topology. Then $e_n$ is continuous,
and the product map
\[
 E:\prod_n A_n\longrightarrow\prod_n F_n',
 \qquad E((a_n)_n)=(e_n(a_n))_n
\]
is a continuous homomorphism: its $n$th coordinate is the
continuous projection followed by $e_n$.
It is injective because equality of its values gives equality
in every coordinate by injectivity of $e_n$.
Both products are compact Hausdorff, so $E$ is a closed embedding.

Apply Lemma~\ref{lem:derived-product-two-generator} to $(F_n)$.
It gives a pro-$p$ group $H$ with a dense generating pair,
hence a topologically finitely generated group, and a closed
embedding $D:\prod_n F_n'\to H$.
The composite of the diagonal embedding, $E$, and $D$ is a
continuous injective homomorphism $P\to H$.
Its source is compact and its target is Hausdorff, so the
compact-source lemma gives the required closed embedding.

The source proves the stronger statement that one two-generator
pro-$p$ group contains every countably based pro-$p$ group.
The statement here allows $H$ to depend on $P$ and records
topological finite generation, which is the conclusion needed
to apply Camina's finitely generated embedding theorem.
\end{proof}
